% Chapitre 6 : Estimation et intervalles de confiance

\section{Pourquoi faire de l'estimation ?}

En pratique, on ne connaît pas la moyenne ou la proportion exacte d'une population. On observe un échantillon et l'on cherche à \textbf{estimer} le paramètre inconnu.

\begin{Ex}
\begin{itemize}
    \item Estimer la résistance moyenne d'une production de béton.
    \item Estimer la proportion de pièces défectueuses sur une chaîne.
    \item Estimer la part de jours défavorables à un type d'intervention.
\end{itemize}
\end{Ex}

\section{Estimateur ponctuel}

\begin{Def}
Un \textbf{estimateur ponctuel} est une formule qui fournit une valeur unique pour approcher un paramètre inconnu.
\end{Def}

\begin{Ex}
\begin{itemize}
    \item La moyenne empirique $\overline{X}$ estime la moyenne $\mu$.
    \item La proportion empirique $\widehat{P}$ estime la proportion $p$.
\end{itemize}
\end{Ex}

\begin{Rmq}
Une estimation ponctuelle est simple à utiliser, mais elle ne donne pas à elle seule la précision de l'estimation.
\end{Rmq}

\section{Intervalle de confiance : idée générale}

\begin{Def}
Un \textbf{intervalle de confiance} est un intervalle aléatoire construit à partir de l'échantillon, censé encadrer le paramètre inconnu avec un niveau de confiance donné.
\end{Def}

\begin{Rmq}
Dans ce cours, on utilisera principalement les niveaux de confiance :
\begin{itemize}
    \item $90\%$ ;
    \item $95\%$ ;
    \item $99\%$.
\end{itemize}
\end{Rmq}

\section{Intervalle de confiance d'une moyenne}

\begin{Prop}
Si la moyenne d'échantillon est approximativement normale, alors un intervalle de confiance asymptotique de niveau $95\%$ pour $\mu$ est
$$
\left[\overline{x}-1.96\frac{\sigma}{\sqrt{n}},\ \overline{x}+1.96\frac{\sigma}{\sqrt{n}}\right]
$$
si $\sigma$ est connu.
\end{Prop}

\begin{Rmq}
Dans de nombreux exercices de ce cours, on remplace $\sigma$ par une valeur fournie par l'énoncé ou par un écart-type estimé lorsque cela est explicitement demandé.
\end{Rmq}

\begin{Ex}
On observe une résistance moyenne de $\overline{x}=32.4$ MPa sur $n=36$ éprouvettes, avec un écart-type connu $\sigma=3$ MPa.

Alors :
$$
IC_{95\%} =
\left[
32.4 - 1.96\times \frac{3}{6},
32.4 + 1.96\times \frac{3}{6}
\right]
= [31.42,\ 33.38].
$$
\end{Ex}

\section{Intervalle de confiance d'une proportion}

\begin{Prop}
Pour un grand échantillon, un intervalle de confiance asymptotique de niveau $95\%$ pour une proportion $p$ est
$$
\left[
\widehat{p} - 1.96\sqrt{\frac{\widehat{p}(1-\widehat{p})}{n}},
\widehat{p} + 1.96\sqrt{\frac{\widehat{p}(1-\widehat{p})}{n}}
\right].
$$
\end{Prop}

\begin{Ex}
Sur $200$ pièces contrôlées, $18$ sont non conformes. Alors
$$
\widehat{p}=\frac{18}{200}=0.09.
$$
Un intervalle de confiance à $95\%$ s'obtient avec
$$
0.09 \pm 1.96 \sqrt{\frac{0.09\times 0.91}{200}}.
$$
\end{Ex}

\section{Marge d'erreur}

\begin{Def}
La \textbf{marge d'erreur} est la demi-largeur de l'intervalle de confiance.
\end{Def}

\begin{Rmq}
La marge d'erreur :
\begin{itemize}
    \item diminue quand $n$ augmente ;
    \item augmente quand le niveau de confiance augmente ;
    \item dépend aussi de la dispersion des données.
\end{itemize}
\end{Rmq}

\section{Comment interpréter un intervalle ?}

\begin{Meth}
Pour commenter un intervalle de confiance :
\begin{enumerate}
    \item on rappelle le paramètre estimé ;
    \item on donne l'intervalle numérique ;
    \item on le compare à une valeur de référence ou à un seuil technique ;
    \item on conclut sans sur-interpréter.
\end{enumerate}
\end{Meth}

\begin{Ex}
Si un seuil de conformité moyen est fixé à $30$ MPa et que l'intervalle à $95\%$ pour la moyenne est $[31.42 ; 33.38]$, alors le seuil $30$ MPa est en dessous de tout l'intervalle. Les données sont donc compatibles avec une moyenne supérieure à $30$ MPa.
\end{Ex}

\section{Ce qu'un intervalle ne dit pas}

\begin{Rmq}
Un intervalle de confiance ne dit pas :
\begin{itemize}
    \item que le paramètre a $95\%$ de chance d'être dans cet intervalle une fois l'intervalle calculé ;
    \item que $95\%$ des observations sont dans l'intervalle ;
    \item qu'il n'existe aucun risque d'erreur de décision.
\end{itemize}
\end{Rmq}

\section{Erreurs fréquentes}

\begin{itemize}
    \item Oublier le facteur $\sqrt{n}$.
    \item Utiliser $p$ à la place de $\widehat{p}$ dans un exercice d'estimation de proportion.
    \item Confondre intervalle sur le paramètre et intervalle contenant les observations.
    \item Donner un résultat numérique sans interprétation.
\end{itemize}

\section{Ce qu'il faut savoir faire}

\begin{itemize}
    \item calculer une estimation ponctuelle de moyenne ou de proportion ;
    \item construire un intervalle de confiance usuel ;
    \item identifier la marge d'erreur ;
    \item conclure par rapport à un seuil de décision.
\end{itemize}
